% introduction_clinical.tex
\subsection{Clinical Structure From Embeddings}

Clinical assessment is these methods' sternest test; two recent studies mark its promise and boundary. \citet{kambeitz-etal-2025} embedded the items of four established psychopathology questionnaires---depression, anxiety, and stress symptoms, schizotypy, autism-related traits, and attenuated psychosis---and compared item-pair embedding similarities with empirical correlations in four datasets ($n$s of 1{,}099 to 39{,}755): the two correlated in every questionnaire ($r = .18$ to $.57$), machine-learning prediction raised this to moderate-to-high, and embedding clusters matched (by adjusted Rand index) both empirical clusters and published subdomains, semantic consistently beating sentiment-based embeddings. \citet{kojima-etal-2026} pressed the most consequential case---whether the general psychopathology factor ($p$-factor) belongs to syndromes or to assessment language---by factor-analyzing 145 symptom items from eight instruments under item-embedding cosines and polychoric correlations from 985 community respondents (matrix correlation $r = .55$): both yielded a dominant general factor with nearly interchangeable loadings (congruence $= .97$), yet no specific factor reached the .90 benchmark. The divergences were informative: responses fused insomnia with somatic complaints---attributed to shared physiology, not shared meaning---while embeddings grouped items merely sharing wording; response-based factor scores predicted lifetime diagnoses better than text-based ones; weighting analyses found the textual information effectively a subset of the response-based. The lesson calibrates the enterprise: instrument language genuinely shadows empirical structure, but lived co-occurrence holds clinically informative variance beyond what wording predicts.
